\begin{afterword}
\summaryhead

The author recalls an episode from an online chatroom. Someone said a friend had chest pains, called an ambulance, and was told by the paramedics that it was nothing; now the pains were worse. The author was puzzled, because he believed paramedics always take anyone with chest pain to hospital. He resolved the puzzle by deciding that the paramedics must have been right, and said so. Later the person reported that the friend had made the story up.

The author concludes that he should have trusted his confusion. The sense that an explanation ``feels a little forced'' is a signal that either your model or the story is wrong. A model is useful for what it rules out, and someone who can explain any outcome knows nothing. Your strength as a rationalist is to be more confused by fiction than by reality. He cites research suggesting that people believe what they hear by default.

\responsehead

The essay turns one remembered anecdote into a named faculty, ``your strength as a rationalist,'' and the anecdote does not support it.

Start with what happened on the page. A stranger asked what to do about a friend whose chest pains were getting worse. The author answered that it ``must really be nothing.'' The essay treats this as a modeling error and an embarrassment. It never says that the advice was dangerous if the story was true, or that the right answer was to call the ambulance again. The whole weight of the lesson falls on the author's failure to detect a lie, none on the risk his answer would have carried.

Then look at the model that produced the confusion. It held that anyone reporting chest pain is ``hauled off by an ambulance instantly,'' that paramedics ``always'' must transport, that emergency rooms must ``treat anyone.'' Each part is false or overstated. Competent adults can refuse transport, and in some US systems up to a quarter of patients do. The law requires emergency screening and stabilization, not treatment of anyone. A reassured patient staying home is an ordinary event. The author's confusion was produced by a false model, and it would have fired just the same on a true story. The essay's own closing sign, ``Either Your Model Is False Or This Story Is Wrong,'' describes his case exactly, with both halves true. Its thesis tells the reader to choose the second.

So the thesis cannot do what it promises. Nobody knows in advance which story is the fiction. A rule that reads confusion as a sign of fiction protects whatever model the reader already holds, and dismisses the anomalies that should have corrected it. The essay's one true general claim, that a model which explains everything explains nothing, is Popper's, and it does not describe the author's mistake: his model forbade the story, and he ignored it.

Finally, the support is thinner than it looks. The resolution of the story rests on the word of the same stranger who told it. The research cited for our tendency to believe by default is a blog post dated the day before and a 1993 study already challenged in 2005, and the reader is not told of the challenge.

In short: an essay that teaches readers to trust their confusion, built on a case where the author's confusion came from a false belief and his advice could have hurt someone.
\end{afterword}
